\documentclass[10pt]{article}

\usepackage[utf8]{inputenc}

\usepackage[T1]{fontenc}

\usepackage{amsmath}

\usepackage{amsfonts}

\usepackage{amssymb}

\usepackage[version=4]{mhchem}

\usepackage{stmaryrd}

\usepackage{graphicx}

\usepackage[export]{adjustbox}

\graphicspath{ {./images/} }



\begin{document}

С. с., в к-рых алгоритмич. тисла образуются сложением узловых, наз. а д д и т и в ны м и. Так, в древнеегипетской (иероглифической) с. с. числа $1,2,3,4$, $5,6,7,8,9,10,19,40$ обозначались соответственно символами



\begin{center}

\includegraphics[max width=\textwidth]{2023_07_09_19d053e368c6c72f3d89g-1}

\end{center}



Те же числа в римской с. с. обозначаются следующим образом: I, II, III, IV, V, VI, VII, VIII, IX, X, XIX, $\mathrm{XL}$. B этой с. с. алгоритмич. числа получаются путем сложения и вычитания узловых. В числительных русского языка отчетливо выражен аддитивно-мультипликативный способ образования алгоритмич. чисел, напр. триста пятьдесят семь.



В нек-рых с. с., наз. а л ф а в и т н ы м и, числа обозначались теми же знаками, что и буквы. Так, древние греки числа от 1 до 9 , а также все десятки и сотни обозначали при помощи последовательных букв алфавита, снабженных черточками. К примеру, числа 803, 833 и 83 записывались так:



\[

\bar{\omega} \bar{\gamma}, \overline{\omega \lambda} \bar{\gamma}, \bar{\pi} \gamma

\]



Алфавитные изображения чисел употребляли славяне и многие другие народы.



С. с. наз. н е по $\mathbf{3}$ и ц и о н н о й, если каждый числовой знак в записи любого числа в ней имеет одно и то же значение. Если же значение числового знака зависит от его расположения в записи числа, то система наз. по $о$ и ц и о н н о й. Римская с. с. - непозиционная. Любое число в вавилонской с. с. записывалось при помощи комбинации двух знаков: вертикального клина и углового клина (пример см. ниже). Эти знаки объединялись в группы от одного до девяти в случае вертикальных клиньев и от одного до пяти в случае угловых клиньев. Вертикальный клин мог обозначать единицу и любую степень числа 60, а угловой клин 10 и произведение 10 на любую степень числа 60. Порядок следования разрядов чисел совпадал с ныне принятым. Так,



\[

7 K K Y Y_{\Delta A d}^{Y Y Y}=1 \cdot 60^{2}+2 \cdot 600+1 \cdot 60+1 \cdot 10+6=4876

\]



В связи с тем, что в вавилонской С. с. отсутствовал знак для пропуска разряда, соответствующий нашему нулю, запись числа не гарантировала однозначность ее прочтения. Точный смысл записи обычно устанавливался из контекста рукописи. Поэтому описанную с. с. принято называть неабсолютно позиционной. Позднее в древнем Вавилоне стали употреблять специальный знак для пропуска разряда. Современная десятичная с. с. - позиционная.



Все известные позиционные с. с. - аддитивно-мультипликативные системы. Позиционный принцип записи чисел в таких системах оправдывается следующей теоремой элементарной теории чисел.



Пусть $q_{0}=1$ и $q_{1}, q_{2}, \ldots$ - последовательность отличных от единицы натуральных чисел. Тогда для любого натурального числа $a$ можно найти одно и только одно натуральное число $n$, для к-рого уравнение



\[

a_{0}+a_{1} q_{1}+a_{2} q_{1} q_{2}+\ldots+a_{n-1} q_{1} \ldots q_{n-1}=a

\]



имеет решение в целых числах $a_{0}, a_{1}, \ldots, a_{n-1}$ таких, पто



$0 \leqslant a_{0}<q_{1}, \ldots, 0 \leqslant a_{n-1}<q_{n-1}, 0<a_{n-1}<q_{n}$.



При әтом уравнению (1) удовлетворяет только один упорядоченный набор (кортеж)



\[

\left\langle a_{n-1}, \ldots, a_{0}\right\rangle

\]



целых чисел с условием (2).



В вавилонской с. с. $q_{1}=10, q_{2}=6, q_{3}=10, q_{6}=6, \ldots$ и т. д. У индейцев племени майя $q_{1}=5, q_{2}=4, q_{3}=18$, $q_{4}=q_{5}=\ldots=20$.



C. с., в к-рой все члены упомянутой в теореме последовательности $q_{1}, \ldots, q_{n}$ равны одному и тому же числу $q$ и в к-рой каждое ив чисел от 0 до $q-1$ обозначается определенным символом, наз. $q-$ и ч н о й 1 с. с. или



\includegraphics[max width=\textwidth, center]{2023_07_09_19d053e368c6c72f3d89g-1(1)}

$q$-ичной с. с. каждое натуральное число обозначается кортежем из указанных символов. Для выполнения сложения и умножения чисел, записанных в $q$-ичной с. с., достаточно иметь таблицы сложения и умножения для всех чисел от 0 до $q-1$.



Лum.: [1] Б а шм а ко в а И. Г., Юшке вич А. Н., Происхождение систем счисления, в кн.: Энциклопедия әлемен-

тарной математики, кн. 1, М.-ЈI., 1951, с. $11-74$; [2] В а н д е p тарной математики, кн. 1, М.-Л., 1951, с. 11-74; [2] В а н д е p $1959 ;$ [3] Ю ш к е в и ч А. П., История математики в ср. века, ка, М., 1961; [5] История математики, т. 1, М., 1970.



СЮРЂЕКЦИЯ, с ю р ъ $е$ т и в н о о В и $\mathrm{e}$ множества $A$ в мотесто $B$ - отображ а такое, что $f(A)=B$. Вместо «f сюръективно» говорят также «f есть отображение множества $A$ на множество $B$ ». С. множества $A$ на $A$ наз. также п е р е с т а н о вк о й множества $A$.



О. А. Иванова.





\end{document}